%
%
\section{Economic Profile and Robustness: Construction and Diagnostics}
\label{app:profile_submission}

This appendix documents the construction and diagnostics behind the
Section~5 economic-profile and robustness exercises. Throughout,
``cobidder'' denotes an \emph{adjudication-anchored exposure label} ---
a firm that co-bid with a firm named in a CADE proceeding --- and never
a finding of participation in any agreement. The profiles below
describe \emph{losing behavior} on the award layer; they are
descriptive inputs to a triage ranking, not proof of conduct.

%--------------------------------------------------------------------
\subsection{Comparison Groups and Profile Variables}
\label{app:profile_groups_submission}

All groups are drawn from the always-loser universe (firms with zero
recorded wins over the calibration window, $N=16{,}843$). The narrow
adjudication-anchored target is the set of $191$ always-loser
cobidders. The $46$ direct CADE defendants are \emph{winners} and are
\emph{excluded} from every cobidder group; the $6$ that appear in the
always-loser universe are excluded from the profile contrasts to keep
the comparison free of named-defendant contamination. The principal
Section~5 contrast is D (FL non-cobidders, $N=2{,}544$) versus E (FL
cobidders, $N=191$). Online Supplement~D lists all group definitions.
% DF-COMPRESS: moved tab:profile_groups_submission to Online Supplement D

Profile variables are computed from award records over the calibration
window and organized in five panels: (A) participation and survival
($T_i$, active years, items per active year); (B) breadth and
concentration (distinct buyers, item groups, modality mix, item-group
and buyer HHIs); (C) defendant-proximity / opportunity layer (observed
contacts $O_i$, opportunity-expected contacts $E_i^{\mathrm{MED}}$,
excess $X_i^{\mathrm{MED}}=O_i-E_i^{\mathrm{MED}}$, cases touched,
contact intensity $O_i/T_i$, share in defendant-bearing cells); (D)
modality-quorum proxies (Convite minimum-quorum share, mean competitors
faced); and (E) timing (first-year and late-half participation shares).
Three variable classes that would sharpen interpretation are
NOT\_OBSERVED in the award layer and used nowhere in Section~5: tender
value / contract size (price is recorded only for the \emph{winner}),
distance-to-winner (requires bid microdata), and fine geography /
specialization. Because of the latter two, the ordinary-loser
alternatives in Section~\ref{sec:cobidder_type} are \emph{narrowed but
not eliminated}.

%--------------------------------------------------------------------
\subsection{Standardized Differences}
\label{app:profile_smd_submission}

Section~5 compares groups~D and~E using standardized mean differences
(Cohen's $d$) rather than $p$-values. With $N$ in the thousands,
$p$-values reject essentially any difference and convey no information
about magnitude; the SMD is scale-free and comparable across
heterogeneous variables. For a continuous variable with group means
$\bar{x}_E,\bar{x}_D$ and standard deviations $s_E,s_D$,
\[
d=\frac{\bar{x}_E-\bar{x}_D}{s_p},
\qquad
s_p=\sqrt{\tfrac{(n_E-1)s_E^2+(n_D-1)s_D^2}{n_E+n_D-2}} ,
\]
with the binary analogue obtained by substituting
$\hat{p}(1-\hat{p})$ for the variances. A positive $d$ means cobidders
(E) are higher. Table~\ref{tab:profile_smd_submission} reports the
top-ranked differences; the full twenty-variable table is in Online
Supplement~D.

\begin{table}[htbp]
\centering
\begin{adjustbox}{max width=\textwidth,center}
\begin{threeparttable}
\caption{Standardized profile differences (top variables),
E (FL cobidders) minus D (FL non-cobidders).}
\label{tab:profile_smd_submission}
\scriptsize
\begin{tabular}{cllr r}
\toprule
Panel & Variable & Label & Raw diff. & Cohen's $d$ \\
\midrule
C & $n\_cases$ & \# CADE cases touched & 1.108 & 3.757 \\
C & $n\_def\_firms$ & \# distinct defendants & 1.335 & 3.117 \\
C & $contact\_intensity$ & $O_i/T_i$ & 0.166 & 1.852 \\
C & $share\_in\_def\_cells$ & part.\ in def.\ cells & 0.210 & 1.762 \\
C & $O_i$ & observed contacts & 7.285 & 1.592 \\
C & $X_i^{\mathrm{MED}}$ & excess contacts & 1.192 & 0.903 \\
A & $n\_years$ & \# active years & 1.324 & 0.841 \\
B & $pregao\_share$ & Preg\~ao share & 0.301 & 0.761 \\
B & $top\_ig\_share$ & top item-group share & 0.122 & 0.516 \\
A & $T_i$ & \# tender-items & 28.557 & 0.505 \\
B & $convite\_share$ & Convite share & $-0.301$ & $-0.761$ \\
\bottomrule
\end{tabular}
\begin{tablenotes}[flushleft]
\scriptsize
\item \textit{Notes:} $n_D=2{,}544$, $n_E=191$ (modality shares use
$n_D=2{,}536$ where modality is defined). Positive $d$: cobidders
are higher. Top eleven of twenty variables by $|d|$; the full table
is in Online Supplement~D. The largest gaps are
defendant-proximity variables ($n\_cases$ $d=3.76$, $n\_def\_firms$
$d=3.12$), which are partly mechanical because the cobidder label is
itself adjudication-anchored. Item-group HHI shows a modest
raw gap ($d=0.45$) that falls to $\approx 0.02$ after opportunity
adjustment.
\end{tablenotes}
\end{threeparttable}
\end{adjustbox}
\end{table}

The defendant-proximity gaps in panel~C are \emph{partly mechanical}:
the cobidder label is defined by adjudication-anchored contact, so
firms with more contact are mechanically more likely to be labeled. We
hold exposure fixed by regressing each profile variable on the cobidder
indicator with controls for $\log E_i^{\mathrm{MED}}$, a binned score
term, breadth, and active years, and by coarsened exact matching on
$E_i^{\mathrm{MED}}$ deciles. Under either adjustment the item-group
HHI gap falls from $d\approx 0.45$ to $\approx 0.02$. Proximity-variable
survival under these adjustments is \emph{partly mechanical} and should
not be read as independent corroboration of conduct.

%--------------------------------------------------------------------
\subsection{Robustness: Monotonicity, Threshold, Markets, and Null Controls}
\label{app:profile_robustness_submission}

Four robustness families support the exposure-ranking reading; full
detail is in Online Supplement~D, which reports the threshold sweep,
the bunching histogram, the market-specific zero-win variants, and the
negative controls. \emph{Monotonicity.} Across participation bins,
cobidder prevalence rises (Spearman $\rho=+0.92$) while mean excess
contact falls ($\rho=-0.93$): the score ranks firms by \emph{exposure},
not collusion intensity. \emph{Threshold.} The $\mathrm{FL14}$ cut
$T_i\geq 14$ equals $\mathrm{median}+1.5\times\mathrm{IQR}$ ($=13.5$,
rounded up) of the always-loser participation distribution; it is an
\emph{administrative} construction, computed ex post and \emph{never
published}. Discrimination changes smoothly across a placebo-threshold
grid with no discontinuity at $14$, and the participation distribution
declines smoothly through the cut (local density ratio $1.06$,
indistinguishable from $1$), so there is no strategic bunching.
\emph{Market-specific zero-win.} Four home-market loser definitions
overlap the global set on only $0.33$--$0.55$ of their losers; the
global definition is not dominated (ROC-AUC $0.939$ vs.\ the best
home-market alternative $0.774$) and is stable to
leave-one-item-group-out re-estimation ($[0.936,0.941]$). \emph{Negative
controls.} Against matched null targets ($B=500$ draws, seed
\texttt{20260603}), the real ROC-AUC ($0.939$) lies far above both null
distributions (means $0.755$ and $0.780$; empirical $p<0.001$). The
placebo does \emph{not} separate on PR-AUC ($p=0.70$) because random
$T_i$-matched anchors touch several times more firms than the $191$ real
positives, inflating the base rate; ROC-AUC, base-rate invariant, is the
decisive metric.
% DF-COMPRESS: moved tab:profile_negcontrol_submission to Online Supplement D

%--------------------------------------------------------------------
\subsection{Limitations}
\label{app:profile_limits_submission}

The defendant-contact variables carry the largest standardized gaps but
are partly a construction artifact; the score ranks firms by exposure to
defendant-adjacent opportunities, not by conduct intensity; ordinary
specialization and geographic-niche alternatives are narrowed but not
eliminated because distance-to-winner, geography, and tender value are
NOT\_OBSERVED. Award-layer profiles describe patterns of \emph{losing};
they are not proof of any agreement, and ``cobidder'' is throughout an
adjudication-anchored exposure label, never an assertion of membership.
